\section{结论展望}

\begin{frame}{研究结论}
  \begin{enumerate}
    \item 构建了面向企业信用风险预警的 \hlb{动态评估框架}。
    \item 实证结果表明：LSTM 在 \hl{AUC、KS、召回率} 上优于 Logistic 回归。
    \item Logistic 回归在 \hl{精确率、可解释性、审计友好性} 上仍有明显优势。
    \item “LSTM 初筛 + Logistic 复核”的混合体系，为本科论文提出了 \hlt{可落地的业务方案}。
  \end{enumerate}
\end{frame}

\begin{frame}{局限与未来展望}
  \begin{columns}[T,onlytextwidth]
    \column{0.48\textwidth}
    \begin{block}{研究局限}
      \begin{itemize}
        \item 主要使用结构化财务指标
        \item 仅比较 LSTM 与 Logistic 两类代表模型
        \item 当前聚焦 1 年期预警任务
      \end{itemize}
    \end{block}

    \column{0.48\textwidth}
    \begin{exampleblock}{未来方向}
      \begin{itemize}
        \item 融合市场、文本等多模态数据
        \item 引入更前沿的时序模型与集成学习
        \item 推进在线学习、滚动预测与可解释性研究
      \end{itemize}
    \end{exampleblock}
  \end{columns}

  \vspace{0.6em}
  \begin{alertblock}{结束语}
    本研究说明：\hl{时序信息确实有价值，但真正可用的风控系统往往来自模型协同，而不是单模型崇拜。}
  \end{alertblock}
\end{frame}
